\documentclass[11pt]{article}
\usepackage[a4paper, top=18mm, bottom=18mm, left=20mm, right=20mm]{geometry}
\usepackage[T2A]{fontenc}
\usepackage[utf8]{inputenc}
\usepackage[ukrainian]{babel}
\usepackage{graphicx}
\usepackage[export]{adjustbox}
\usepackage{amsmath}
\usepackage{amssymb}
\graphicspath{{/app/Quiz/Scripts/2023/ProgAll/15BinTree/}{/app/Quiz/Scripts/2020/Mod2020_4/BinTree/}{/app/Quiz/Scripts/2021/Mod2021_3/BinTree/}}
\setlength{\parindent}{0pt}
\setlength{\parskip}{4pt}
\pagestyle{empty}
\begin{document}
%begin
{\large\textbf{Контрольна робота}}

\textbf{Варіант \No\,1}\hfill \textit{ПІБ:}\,\underline{\hspace{60mm}}
\vspace{4mm}

%beginitem
1. 

try:
    print(6,end="")
    print(3j<6j,end="")
    print(1,end="")
except Exception: print(2,end="")
except TypeError: print(5,end="")
else: print(4,end="")
finally: print(9,end="")

%enditem
%beginitem
2. %goodpath:Scripts/2018/Mod2018_1/03-good.txt
Відмітити все, що є коректним оператором С++:
( 1 ) double(2)/87

( 2 ) if ('x'><'y') a=6;

( 3 ) cin>>x;

( 4 ) if a<b a=0;

%enditem
%beginitem
3. Що буде виведено за виконання фрагменту коду?

%code
print('R', end=' ')
i=12
while i>6:
    i+=-3
print(i)
%code


%enditem
%beginitem
4. 
try:
    res = not8<8.0 and not5.0>=5 and True<=3
    if isinstance(res, bool):
        print(f'bool;{res}')
    elif isinstance(res, int):
        print(f'int;{res}')
    elif isinstance(res, float):
        print(f'float;{res:.2f}')
    else:
        print('???')
except:
    print('Z')

%enditem
%beginitem
5. 
Змінні $next$ та $prev$ вузла списку позначають наступний та попередній за ним вузол відповідно. Починаючи з голови, у список записано послідовні цілі числа від $-38$ до $47$. Нехай змінна $p$ позначає вузол, що зберігає значення $-2$.
Яке значення зберігається у вузлі $p.prev.prev.prev.prev.prev$ ?

%enditem
%beginitem
6. 
g=[10,11,12,13,14]; print(g.pop(1),end=' '); print(g)
%enditem
%beginitem
7. 
Написати програму, що запитує у користувача значення дійсних параметрів $x$ та $y$, обчислює для них значення виразу 
$$\frac{\log_{2} x + 54}{(\tg x + 0.8) (y + 70)}$$ 
та повiдомляє введенi данi та результат обчислення.
 Оцінюється правильність та якість коду.


%enditem
%beginitem
8. 
Записати ВИРАЗ, значенням якого є словник, ключами якого є цілі числа від 17 до 2007 (включно), що не діляться на жодне з чисел 2, 3, а ключам відповідають їх квадратні корені 
%enditem
%beginitem
9. Що буде виведено за виконання фрагменту коду?

%code
#include <iostream>
using namespace std;

int h(int a){
    int u = 68;
    if (a < 2) 
        return 6;
    if (a >= -2)
         u = 3;
    else
         u = 1;
    return u;
}

int main(){
    cout << h(-2);
    return 0;
}
%code


%enditem
%beginitem
10. Що буде виведено за виконання фрагменту коду?

%code
print('R', end=' ')
g=('a','b','c','d','e',0,1,2,3)
print(g[2])
%code


%enditem
%beginitem
11. Що буде виведено за виконання фрагменту коду?

%code
9**0.5*-2
%code

%enditem
%beginitem
12. 
i=12
   while i>6: i+=-3
   print(i)
%enditem
%beginitem
13. 
Знайти кількість відношень на n-елементній множині, які нерефлексивні та несиметричні
%enditem
%beginitem
14. Що буде виведено за виконання фрагменту коду?

%code
a = 6
b = 3
c = 1

def h(a):
global c    a=3
    b*=4
    c=1
    return a+b+c

a, b, c = 2, 5, 4
try:
    print(h(a), a, b, c, sep=';')
except:
    print('Z', a, b, c, sep=';')
%code

%enditem
%beginitem
15. 
x,y,z=0,3,1;  x,x,y,z=6,z,y,x;  print(x,y,z)
%enditem
%beginitem
16. 

print('R', end=' ')
i=12
while i>6:
    i+=-3
print(i)


%enditem
%beginitem
17. 

def h(a):
    u=68
    if a<2: 
        return 6
    if a>=-2:
         u=3
    else:
         u=1
    return u

print(h(-2))

%enditem
%beginitem
18. 
Змінні \kw{next} та \kw{prev} вузла списку позначають наступний та попередній за ним вузол відповідно. Починаючи з голови, у список записано послідовні цілі числа від~$-38$ до~$47$. Нехай змінна \kw{p} позначає вузол, що зберігає значення~$-2$.
Яке число зберігається у вузлі \kw{p.prev.prev.prev.prev.prev} ?

%enditem
%beginitem
19. Що буде виведено за виконання фрагменту коду?

%code
print(8 < "4")
%code

%enditem
%beginitem
20. Що буде виведено за виконання фрагменту коду?

%code
#include <iostream>
using namespace std;

bool f(int n){
    cout<<"f";
    return n;
}

int main(){
    cout<<(f(-2) and f(-6));
    return 0;
}
%code

%enditem
%beginitem
21. 
(Задача за частиною Пашка А.О. Наводиться в авторській редакції.)\par {\hspace {0.5 cm}}³дбудеться деяка подія $W$ чи ні залежить від двох факторів $A$ та $X$. За результатами статистичного експерименту (100 раз для кожного значення фактора $A$,  200 раз для кожного значення фактора $X$) подія $W$ відбулася:\par {\hspace {0.5 cm}}$(n+17)$ раз для значення $A(a)$, $(n+21)$ раз для значення $A(b)$, $(n+50)$ раз для значення $A(c)$, $(n+57)$ раз для значення $A(d)$, та відповідно $(2n+100)$ раз для значення $X(x)$, $(4n+17)$ раз для значення $X(y)$, $(n+150)$ раз для значення $X(z)$.\par
{\hspace {0.5 cm}}Знайти ймовірність того що подія $W$ не відбулася (умовна ймовірність) при значеннях факторів $A(b)$, $X(y)$. Фактори є незалежними, $n$ – номер цього екзаменаційного білета.\par
%enditem
%beginitem
22. 
try:
    res = 5/5*5*5
    if isinstance(res, bool):
        print(f'bool;{res}')
    elif isinstance(res, int):
        print(f'int;{res}')
    elif isinstance(res, float):
        print(f'float;{res:.2f}')
    else:
        print('???')
except:
    print('Z')

%enditem
%beginitem
23. 
try:
    for a in range(-8, -8, -3):
        if a<3:
            break
        print(a, end=';')
        if a<3:
            break
    else:
        print('end',end=';')
    print(a)
except:
    print('ERR')

%enditem
%beginitem
24. 

\begin{center}
	\adjustimage{max size={0.8\linewidth}{0.8\paperheight}}
	{../BinTree/trees_exp/140.png}
\end{center}
Указати послідовність міток вершин дерева, зображеного на рис., що утвориться в результаті обходу дерева в прямому порядку. Мітки записувати через пробіл.\par Обхід дерева починається з кореня (на рис. це верхня вершина).
%answer:140 прямому

%enditem
%beginitem
25. 
Значенням змінної \kw{a} є цілочисловий масив типу $numpy.ndarray$ розмірів $5{\times}5$. На скільки збільшиться сума елементів масиву \kw{a} після виконання
\begin{verbatim}
a[-3] += 3
\end{verbatim}


Якщо код некоректний, то в якості відповіді зазначити ERROR.



%enditem
%beginitem
26. Що буде виведено за виконання фрагменту коду?

%code
print('R', end=' ')
g=[0,1,2,3,4,5,6,7,8,9]
print(g[:-8:-3])
%code


%enditem
%beginitem
27. Що буде виведено за виконання фрагменту коду?

%code
print(8<"4")
%code

%enditem
%beginitem
28. Що буде виведено за виконання фрагменту коду?

%code
for a in range(-3,-3+4,-3):
    if a<3:
        continue
    print(a, end=' ')
    a=-3
    if a<3:
        break
else:
    print(13, end=' ')
print(a, end=' ')
%code

%enditem
%beginitem
29. 
В умовах попередньої задачі було виконано p->prev=p->next->next;

Чому дорівнює значення p->prev->prev->prev->prev->n ?
%enditem
%beginitem
30. Що буде виведено за виконання фрагменту коду?
%code
if (6 < 6)
    cout << "g";
else
    cout << "g";
%code


%enditem
%beginitem
31. Що буде виведено за виконання фрагменту коду?

%code
cout << 12 % 12 % 12 % 12;
%code

%enditem
%beginitem
32. Що буде виведено за виконання фрагменту коду?

%code
cout << 10 % 10 % 10 % 10;
%code

%enditem
%beginitem
33. Що буде виведено за виконання фрагменту коду?

%code
print('R', end=' ')
g = ('a','b','c','d','e',0,1,2,3)
print(g[2])
%code


%enditem
%beginitem
34. Що буде виведено за виконання фрагменту коду?
%Оригинал был утерян. Требует отладки!
%code
_a = 0

class A:
    _a = 1
    
    def __init__(self):
        self._a = 2
        _a = 3
        A._a = 4

obj = A()
obj._a = 5
try:
    print(_a, end=' ')
    print(obj._a, end=' ')
    print(A._a, end=' ')
    print(obj._a, end=' ')
    print(obj._A__a)
except:
    print('ERR')
%code

%enditem
%beginitem
35. Що буде виведено за виконання фрагменту коду?
%code

x,y,z=0,3,1
x,x,y,z=6,z,y,x
print(x,y,z)
%code

%enditem
%beginitem
36. 

a,b,c=6,3,1
def h(a):
global c    a=3
    b*=4
    c=1
    return a+b+c

a,b,c=2,5,4
print(h(a),a,b,c)

%enditem
%beginitem
37. Що буде виведено за виконання фрагменту коду?

%code
print('R', end=' ')
g = [0,1,2,3,4,5,6,7,8,9]
print(g[:-8:-3])
%code


%enditem
%beginitem
38. 
for a in range(-8,-8,-3):
    if a<3:
        break
        print(a,end=' ')
    if a<3:
        break
else:
    print('end', end=' ')
print(a, end=' ')

%enditem
%beginitem
39. 
Записати ВИРАЗ, значенням якого є список, що складається з кубівцілих чисел від 17 до 2007 (включно), що не діляться на жодне з чисел 2, 3, записаних в оберненому порядку.
%enditem
%beginitem
40. 
Записати ВИРАЗ, значенням якого є список, що складається з кубівцілих чисел від 17 до 2007 (включно), що не діляться на жодне з чисел 2, 3, записаних в оберненому порядку.
%enditem
%beginitem
41. Що буде виведено за виконання фрагменту коду?

%code
#include <iostream>

int h(int a, int b){
    int c = 68;
    if (a < 2)
        return 6;
    if (b >= -2)
         c = 3;
    else 
        c = 1;
    return c;
}

int main(){
    std::cout << h(-2, -2);
    return 0;
}
%code

%enditem
%beginitem
42. 
a,b,c=6,3,1
    def h(a):
      global c      a=3
      b=4
      c=1
      return a+b+c
    a,b,c=2,5,4
    print(h(a),a,b,c)
%enditem
%beginitem
43. %goodpath:Scripts/2019/Mod2019_1/Lex/03-good.txt
Відмітити все, що є коректним оператором С++:
( 1 ) x-=2

( 2 ) x in not y

( 3 ) x is y is z

( 4 ) y=/x

%enditem
%beginitem
44. Що буде виведено за виконання фрагменту коду?

%code
cout << (!8<8.0 and !5.0>= 5 and true<=3);
%code

%enditem
%beginitem
45. 

\begin{center}
	\adjustimage{max size={0.8\linewidth}{0.8\paperheight}}
	{../BinTree/trees_exp/140.png}
\end{center}
Указати послідовність міток вершин дерева, зображеного на рис., що утвориться в результаті обходу дерева в прямому порядку. Мітки записувати через пробіл.\par Алгоритм починає роботу з кореня (на рис. це верхня вершина).
%answer:140 прямому

%enditem
%beginitem
46. Що буде виведено за виконання фрагменту коду?

%code
print(8>8.0>5.0>5)
%code

%enditem
%beginitem
47. 
a = 6
b = 3
c = 1

def h(a):
global c    a = 3
    b = 4
    c = 1
    return a + b + c

a = 2
b = 5
c = 4

try:
    print(h(a), a, b, c, sep=';')
except:
    print('Z', a, b, c, sep=';')

%enditem
%beginitem
48. Що буде виведено за виконання фрагменту коду?

%code
try:
    i = 12
    while i>6:
        i += -3
    print(i, end=';')
except:
    print('Z')

%code

%enditem
%beginitem
49. Що буде виведено за виконання фрагменту коду?

%code
print(not8<8.0 and not5.0>=5 and True<=3)
%code

%enditem
%beginitem
50. 
8<"4"
%enditem
%beginitem
51. Що буде виведено за виконання фрагменту коду?

%code
def h(a,b):
    c=68
    if a<2:
        return 6
    if b>=-2:
         c=3
    else: 
        c=1
    return c

print(h(-2,-2))
%code

%enditem
%beginitem
52. Що буде виведено за виконання фрагменту коду?

%code
print(5/5*5*5)
%code

%enditem
%beginitem
53. 
a = 6
b = 3
c = 1

def h(a):
global c    a = 3
    b = 4
    c = 1
    return a + b + c

a = 2
b = 5
c = 4

try:
    print(h(a), a, b, c, sep=';')
except:
    print('ERR', a, b, c, sep=';')

%enditem
%beginitem
54. 
g=[10,11,12,13,14]; del g[-3:]; print(g)
%enditem
%beginitem
55. 
a,b,c=6,3,1
    def h(a):
      global c      a=3
      b*=4
      c=1
      return a+b+c
    a,b,c=2,5,4
    print(h(a),a,b,c)
%enditem
%beginitem
56. Що буде виведено за виконання фрагменту коду?

%code
g=[10,11,12,13]
print(g.pop(2), end=' ')
print(g)
%code


%enditem
%beginitem
57. 
g=[0,1,2,3,4,5,6,7,8,9]; print(g[:-8:-3])
%enditem
%beginitem
58. %goodpath:Scripts/2019/Mod2019_1/Lex/01-good.txt
Відмітити все, що є однією правильною лексемою:
( 1 ) 0.2E-3

( 2 ) 0x17

( 3 ) false

( 4 ) 'a'a'

%enditem
%beginitem
59. 
for a in range(-8,-8,-3):
    if a<3:
        break
        print(a,end=' ')
    if a<3:
        break
    else:
        print('end',end=' ')
    print(a)

%enditem
%beginitem
60. 
Рядки журналу через двокрапку містять ім'я користувача (ідентифікатор), час події,тип події, код події, наприклад
\begin{verbatim}
s="username : 2022-05-05 13:26 : ERROR : 404"
\end{verbatim}
у коді 
\begin{verbatim}
res = re.fullmatch('???', s) 
s = ''
code_str = ??? 
\end{verbatim}
чим треба замінити кожен з ???, щоб значенням code\_str був код події? 



%enditem
%beginitem
61. Що буде виведено за виконання фрагменту коду?

%code
#include <iostream>
using namespace std;

int h(int a, int b){
    int c = 68;
    if (a < 2)
        return 6;
    if (b >= -2)
         c = 3;
    else 
        c = 1;
    return c;
}

int main(){
    cout << h(-2, -2);
    return 0;
}
%code

%enditem
%beginitem
62. 
def f(arg, n):
    n = 6
    tmp = arg
    del tmp[-3:]
    return len(tmp)>3

try:
    g = [10,11,12,13,14]
    n = 3
    f(g, n)
    print(n, end=';')
    for el in g:
        print(el, end=';')
except:
    print('ERR')

%enditem
%beginitem
63. %goodpath:Scripts/2018/Mod2018_1/01-good.txt
Відмітити все, що є однією правильною лексемою:
( 1 ) 0.5

( 2 ) 2*6

( 3 ) While

( 4 ) x^2

%enditem
%beginitem
64. 
Записати ВИРАЗ, значенням якого є словник, ключами якого є цілі числа від 17 до 2007 (включно), що не діляться на жодне з чисел 2, 3, а ключам відповідають їх квадратні корені 
%enditem
%beginitem
65. 
def f(arg, n):
    n = 6
    tmp = arg
    del tmp[-3:]
    return len(tmp)>3

try:
    g = [10,11,12,13,14]
    n = 3
    f(g, n)
    print(n, end=';')
    for el in g:
        print(el, end=';')
except:
    print('Z')

%enditem
%beginitem
66. 
i=6
   while i<12: i+=1
   print(i)
%enditem
%beginitem
67. Що буде виведено за виконання фрагменту коду?

%code
def h(a):
    u=68
    if a<2: 
        return 6
    if a>=-2:
         u=3
    else:
         u=1
    return u

print(h(-2))
%code

%enditem
%beginitem
68. 

print('R', end=' ')
g=[10,11,12,13,14]
del g[-3:]
print(g)


%enditem
%beginitem
69. Що буде виведено за виконання фрагменту коду?

%code
print('R', end=' ')
i=6
while i<12:
    i+=1
print(i)
%code


%enditem
%beginitem
70. Що буде виведено за виконання фрагменту коду?
code
_a = 0

class A:
    _a = 1
    
    def __init__(self):
        self._a = 2
        _a = 3
        A._a = 4

obj = A()
try:
    print(_a, end=' ')
    print(obj._a, end=' ')
    print(A._a, end=' ')
    print(obj._a, end=' ')
    print(obj._A__a, end=' ')
except:
    print('ERR')
code

%enditem
%beginitem
71. 

a,b,c=6,9,7
def h(a,b,c):
    print(a,b,c,end="")

h(b=0,c=3,1)
print(a,b,c)


%enditem
%beginitem
72. Що буде виведено за виконання фрагменту коду?

%code
cout << (!8 < 8.0 and !5.0 >= 5 and true <= 3);
%code

%enditem
%beginitem
73. 
a = 6
b = 3
c = 1

def h(a):
global c    a = 3
    b *= 4
    c = 1
    return a + b + c

a = 2
b = 5
c = 4

try:
    print(h(a), a, b, c, sep=';')
except:
    print('ERR', a, b, c, sep=';')

%enditem
%beginitem
74. Що буде виведено за виконання фрагменту коду?

%code
g = [0,1,2,3,4,5,6,7,8,9]
print(g[:-8:-3])
%code


%enditem
%beginitem
75. Що буде виведено за виконання фрагменту коду?

%code
8<"4"
%code

%enditem
%beginitem
76. 
def f(n): print("f", end=""); return n
   print(f(-2) and f(-6))
%enditem
%beginitem
77. 

\begin{center}
	\adjustimage{max size={0.8\linewidth}{0.8\paperheight}}
	{../15BinTree/trees_exp/140.png}
\end{center}
Указати послідовність міток вершин дерева, зображеного на рис., що утвориться в результаті обходу дерева в прямому порядку. Мітки записувати через пробіл.\par Алгоритм починає роботу з кореня (на рис. це верхня вершина).
%answer:140 прямому

%enditem
%beginitem
78. Що буде виведено за виконання фрагменту коду?

%code
#include <iostream>
using namespace std;

int a = 6, b = 3, c = 1;

int h(int &a){
int c;    a += 3;
    b = 4;
    c = 1;
    return a + b + c;
}

int main(){
    inta = 2;
    intb = 5;
    intc = 4;
    cout << h(a) << ':';
    cout << a << ':' << b << ':' << c;
    return 0; 
}
%code

%enditem
%beginitem
79. 
Записати ВИРАЗ, значенням якого є множина, що складається з кубівцілих чисел від 17 до 2007 (включно), що не діляться на жодне з чисел 2, 3.
%enditem
%beginitem
80. Що буде виведено за виконання фрагменту коду?

%code
i=6
while i<12:
    i+=1
print(i, end=' ')
%code


%enditem
%beginitem
81. Що буде виведено за виконання фрагменту коду?
%code
#include <iostream>
using namespace std;

int f(int &x, int &y){
    x = 8;
    y+= 4;
    return x;
}

int main(){
    int a = 1, b = 2;
    a = f(a, a);
    cout << a << ":" << b <<':';
    {
        int a = 5, b = 7;
        cout << ((a<5) && ((b+=1) < 5)) << ':';
        cout << a << ":" << b << ':';
    }
    {
        inta = 6;
        cout << a << ':';
    }
    cout << a << ":" << b << endl;
    return 0;
}
%code


%enditem
%beginitem
82. 
Скільки 2017-елементних підмножин множини натуральних від 1 до 20172017 містять принаймні одне число, що більше 172006
%enditem
%beginitem
83. 
try:
    res = 8<"4j"
    if isinstance(res, bool):
        print(f'bool;{res}')
    elif isinstance(res, int):
        print(f'int;{res}')
    elif isinstance(res, float):
        print(f'float;{res:.2f}')
    else:
        print('???')
except:
    print('ERR')

%enditem
%beginitem
84. 

print('R', end=' ')
g=[10,11,12,13,14]
print(g.pop(1), end=' ')
print(g)


%enditem
%beginitem
85. Що буде виведено за виконання фрагменту коду?

%code
cout << (8 < 8.0 >= 5.0 <= 5);
%code

%enditem
%beginitem
86. 
Написати програму, що запитує у користувача значення дійсних параметрів $x$ та $y$, обчислює для них значення виразу 
$$\frac{\log_{2} x + 54}{(\tg x + 0.8) (y + 70)}$$ 
та повiдомляє введенi данi та результат обчислення.


%enditem
%beginitem
87. 
not8<8.0 and not5.0>=5 and True<=3
%enditem
%beginitem
88. 

print('R', end=' ')
for a in range(-8,-8,-3):
    if a<3:
        break
        print(a, end=' ')
    if a<3:
        break
    else:
        print('end', end=' ')
    print(a)

%enditem
%beginitem
89. Що буде виведено за виконання фрагменту коду?
%code
for a in range(-8,-8,-3):
    if a<3:
        break
        print(a,end=' ')
    if a<3:
        break
else:
    print('end', end=' ')
print(a, end=' ')
%code

%enditem
%beginitem
90. Що буде виведено за виконання фрагменту коду?
code
a = 0
k = 1

class A:
    a = 2
    
    def __init__(self):
global a        self.a = 3
        a = 4
        A.a = 5

obj = A()
print(a, k, A.a, obj.a, sep=' ')
code

%enditem
%beginitem
91. Що буде виведено за виконання фрагменту коду?

%code
5/5*5*5
%code

%enditem
%beginitem
92. 
Проект складається з од���ї одиниці трансляції 1.cpp, яка починається таким фрагментом\par
long h(){ return my134();}\par
Функція my134(); означена в 2.cpp.\par
Що треба зробити для того, щоб 1.cpp успішно відкомпілювався?\par\vspace{1.5cm}
Що треба зробити для того, щоб за проектом зібрався виконуваний код?\par

%enditem
%beginitem
93. Що буде виведено за виконання фрагменту коду?

%code
print('R', end=' ')
g=[10,11,12,13,14]
print(g.pop(1), end=' ')
print(g)
%code


%enditem
%beginitem
94. 
$a$ є цілочисловим масивом типу $numpy.ndarray$ розмірів $2{\times}2$. Сума елементів масиву $a$ дорівнює 27. Чому дорівнює сума елементів масиву $a$ після виконання 
\begin{verbatim}
a += 27
\end{verbatim}


Якщо код некоректний, то в якості відповіді зазначити ERROR.



%enditem
%beginitem
95. 
try:
    i = 12
    while i>6:
        i += -3
    print(i, end=';')
except:
    print('Z')

%enditem
%beginitem
96. Що буде виведено за виконання фрагменту коду?

%code
a,b,c=6,3,1
def h(a):
global c    a=3
    b=4
    c=1
    return a+b+c

a,b,c=2,5,4
print(h(a),a,b,c)
%code

%enditem
%beginitem
97. 
У папці XX розташовано проект, до якого входить синаксично правильна одиниця трансляції 1.cpp з рядком\par
{\#}include "2.h"
\parФайла 2.h у папці XX нема, але він є в папці YZ. Що треба зробити, щоб одиниця трансляції 1.cpp, могла успішно відкомпілюватися?Переміщувати, копіювати та змінювати файли з кодом заборонено.\par\vspace{1.5cm}
%enditem
%beginitem
98. 
try:
    res = 9**0.5*-2
    if isinstance(res, bool):
        print(f'bool;{res}')
    elif isinstance(res, int):
        print(f'int;{res}')
    elif isinstance(res, float):
        print(f'float;{res:.2f}')
    else:
        print('???')
except:
    print('ERR')

%enditem
%beginitem
99. Що буде виведено за виконання фрагменту коду?

%code
print('R', end=' ')
g=[10,11,12,13,14]
del g[-3:]
print(g)
%code


%enditem
%beginitem
100. 
try:
    i = 12
    while i>6:
        i += -3
    print(i, end=';')
except:
    print('ERR')

%enditem




































































































%end
\newpage
\end{document}


